\chapter{System organization}\label{ch:system-organization}

\whatyoufind{The four principles the simulator is built on and the two base classes that
carry them; the clock model, which is the reason every latency it reports is comparable; the
map of components and the classes behind them; and the exact system that
\file{MultiCoreSystem.csv} builds when you run the shipped configuration --- the system every
later chapter's figures are drawn on.}

\section{Design principles}

Octopus models a multicore memory hierarchy as a \emph{network of clocked, configurable
components} joined by bounded, back-pressured \emph{interfaces}. A memory request is a
first-class \code{Message} object that travels this network; a single global clock steps every
component; and every hop is timestamped for the Logger. Four principles shape the code, and
they are realised through two base classes that every component inherits
(\cref{tab:base-classes}).

\begin{itemize}
\item \textbf{Modularity.} Each component is self-contained behind a narrow interface, so
  changing one does not ripple into others. A cache controller does not know whether it is an
  L1 or the LLC, nor which protocol it runs; a bus does not know what it carries.
\item \textbf{Extensibility.} Inheritance and polymorphism: a new arbiter, replacement policy,
  topology or memory model is a subclass with a handful of virtual methods; a new coherence
  protocol is a CSV file. The effect shows in the code size of protocols derived from a common
  family (\cref{ch:coherence-engine}).
\item \textbf{Configurability.} Every component derives from \code{Configurable}: it reads its
  defaults from its own CSV, a parent may override a child's, and the command line overrides
  everything, so the whole design space is reachable without recompiling
  (\cref{ch:configuration}).
\item \textbf{Integrability.} The two endpoints --- the request source and the memory sink ---
  are swappable for external simulators: a trace-driven CPU or gem5 on one side, a fixed-latency
  model or MCsim's DDR controllers on the other, with the hierarchy between them unchanged
  (\cref{ch:integration-modes}).
\end{itemize}

\begin{table}[htbp]\centering\small
\caption{The two base classes every component inherits.}\label{tab:base-classes}
\begin{tabularx}{\textwidth}{@{}lX@{}}
\toprule base class & gives every component \\ \midrule
\code{ClockedObj} & \code{m\_clk\_period} and \code{cycleProcess()}: the component can be advanced one tick by the \code{ClockManager} \\
\code{Configurable} & \code{name}, \code{parent\_name}, \code{parameters}, \code{parseFile()}: hierarchical parameters read from CSV and overridable from the CLI \\
\bottomrule
\end{tabularx}
\end{table}

\section{The clock model}

\begin{figure}[htbp]\centering
% The clock model: one ClockManager advances every ClockedObj by one tick; components
% never call each other, they only exchange Messages through interfaces between ticks.
\begin{tikzpicture}[
  font=\sffamily\small, >=Latex,
  box/.style={draw=octrule, fill=white, rounded corners=2pt, minimum height=9mm, inner sep=5pt, align=center},
  clk/.style={box, fill=octctl, draw=octamber, minimum width=34mm},
  comp/.style={box, fill=octmem, draw=octrule, minimum width=36mm},
  iface/.style={box, fill=octio, draw=octpurple, minimum width=22mm, minimum height=6mm, font=\sffamily\scriptsize},
  wire/.style={->, thick, octink},
  msg/.style={->, thick, octpurple, dashed},
]
  \node[clk] (cm) {\textbf{ClockManager}\\\footnotesize \code{init()} $\rightarrow$ \code{run()}: one \code{clkStep} per tick};
  \node[comp, below left=11mm and 18mm of cm] (cpu) {\textbf{CPU}\\\footnotesize \code{cycleProcess()}};
  \node[comp, right=6mm of cpu] (l1) {\textbf{CacheController}\\\footnotesize \code{cycleProcess()}};
  \node[comp, right=6mm of l1] (ic) {\textbf{Interconnect}\\\footnotesize \code{cycleProcess()}};
  \node[comp, right=6mm of ic] (mem) {\textbf{MainMemory}\\\footnotesize \code{cycleProcess()}};
  \foreach \c in {cpu,l1,ic,mem} \draw[wire] (cm.south) -- ++(0,-4mm) -| (\c.north);
  \node[iface, below=7mm of $(cpu.south)!0.5!(l1.south)$] (i1) {CommunicationInterface};
  \node[iface, below=7mm of $(l1.south)!0.5!(ic.south)$] (i2) {CommunicationInterface};
  \node[iface, below=7mm of $(ic.south)!0.5!(mem.south)$] (i3) {CommunicationInterface};
  \draw[msg] (cpu.south) |- (i1.west);  \draw[msg] (i1.east) -| (l1.south);
  \draw[msg] (l1.south) |- (i2.west);   \draw[msg] (i2.east) -| (ic.south);
  \draw[msg] (ic.south) |- (i3.west);   \draw[msg] (i3.east) -| (mem.south);
  \node[below=3mm of i2, font=\sffamily\footnotesize, text=octquiet, align=center] {messages are pushed into an interface on one tick and drained by the next component on \emph{its} tick:\\no component ever calls another synchronously, and every interface is bounded and back-pressured};
  \node[right=2mm of cm, font=\sffamily\footnotesize, text=octquiet, align=left] {every component is a \code{ClockedObj}\\with its own \code{m\_clk\_period}};
\end{tikzpicture}

\caption{The clock model. One \code{ClockManager} advances every \code{ClockedObj} by one tick;
components exchange messages only through interfaces, which the next component drains on its
own tick.}
\label{fig:clock-loop}
\end{figure}

The \code{ClockManager} owns the simulation loop. On each \code{clkStep} it advances every
\code{ClockedObj} whose period has elapsed by calling its \code{cycleProcess()}
(\cref{fig:clock-loop}). Components never call each other synchronously: they communicate only by
placing \code{Message}s into interfaces, which the receiving component drains when it is next
stepped. That decoupling is what makes the hierarchy modular and back-pressured --- an interface
that is full simply refuses the push, and the sender tries again on its next tick.

Because all components step under one manager, latencies are consistent system-wide: the Logger
measures every stage of a request in the \emph{originating core's} clock, whichever component
stamps it (\cref{ch:monitoring}). Each component carries its own \code{m\_clk\_period}, so a
memory model may legitimately tick at a different rate from the cores; in the shipped
configuration every period is 100 and one tick is one cycle everywhere.

\begin{notebox}[Design note: a clock-period bug worth knowing about]
The MCsim bridge once ticked the DRAM model once per \emph{manager} step rather than once per
LLC cycle --- a hundred DRAM cycles per core cycle --- so DRAM looked like a one-cycle device.
The LLC's period is now passed into the bridge; any DRAM number produced before that fix is
untrustworthy. The lesson generalises: a component with its own period must be stepped by the
manager, never by whoever happens to call it.
\end{notebox}

\section{Component map}

\figmed{SystemConfiguration3.pdf}{What a configured system can contain: clusters of cores with
private caches on direct interconnects, further cache levels inside or across clusters, a
shared last-level cache and main memory --- all of it chosen by configuration, not
hard-coded.}{fig:organization}{0.62}

A system is clusters of CPU cores, each with private caches and a direct interconnect, joined
through configurable interconnects to a shared LLC and main memory (\cref{fig:organization}).
\Cref{tab:roles} names the class behind each role; \cref{fig:classes} shows how they relate.

\begin{table}[htbp]\centering\small
\caption{Roles and the classes that play them.}\label{tab:roles}
\begin{tabularx}{\textwidth}{@{}lp{0.47\textwidth}X@{}}
\toprule role & class(es) & notes \\ \midrule
clock & \code{ClockManager}, \code{ClockedObj} & drives \code{cycleProcess()} each tick \\
configuration & \code{Configurable} & hierarchical CSV + CLI parameters \\
core & \code{CPU} (trace), \code{ExternalCPU} (gem5) & issues \code{Message}s; the trace CPU has an out-of-order window of \key{m\_number\_of\_OoO\_requests} \\
transport & \code{CommunicationInterface} & \code{pushMessage} / \code{peekMessage} / \code{popFrontMessage}; bounded \\
message & \code{Message} & \code{msg\_id}, \code{addr}, \code{cycle}, \code{owner}, \code{source}, \code{data}, \code{kind} \\
cache & \code{BaseController} $\rightarrow$ \code{CacheController} $\rightarrow$ \code{CacheControllerExclusive}, \code{CacheController\_End2End}, \code{CacheControllerDirectory} & processing queue + handlers (\cref{ch:cache-controller}) \\
coherence & \code{CoherenceProtocolHandler} $\rightarrow$ MSI / MESI / MOESI, their LLC and directory variants & driven by \code{FSMReader} (\cref{ch:coherence-engine}) \\
data & \code{CacheDataHandler} $\rightarrow$ \code{CacheDataHandler\_COTS} & array, MSHRs, write-back buffer, replacement, port arbiter \\
interconnect & \code{InterconnectTopology} + \code{InterconnectController} + \code{Arbiter} & bus, mesh, NoC (\cref{ch:interconnects}) \\
memory & \code{MainMemoryController} (fixed latency), \code{MCsimInterface} & the sink (\cref{ch:main-memory}) \\
monitoring & \code{Logger}, \code{DebugPrint} & \cref{ch:monitoring} \\
\bottomrule
\end{tabularx}
\end{table}

\figwide{Octopus.pdf}{The class diagram. Cache-related classes in red, interconnect-related in
blue, utilities in white.}{fig:classes}

\section{The system as shipped}\label{sec:shipped-system}

\begin{figure}[htbp]\centering
% The system MultiCoreSystem.csv builds: four cores with private L1s on a TripleBus, one
% shared inclusive LLC, a point-to-point memory bus and a fixed-latency main memory.
\begin{tikzpicture}[
  font=\sffamily\small, >=Latex,
  box/.style={draw=octrule, rounded corners=2pt, inner sep=4pt, align=center, minimum height=8mm},
  core/.style={box, fill=white, minimum width=27mm, minimum height=19mm},
  cpu/.style={box, fill=octctl, draw=octamber, minimum width=23mm, minimum height=6mm, font=\sffamily\scriptsize},
  l1/.style={box, fill=octmem, minimum width=23mm, minimum height=6mm, font=\sffamily\scriptsize},
  bus/.style={box, fill=octio, draw=octpurple, minimum height=11mm},
  llc/.style={box, fill=octmem, draw=octrule, minimum height=12mm},
  mem/.style={box, fill=white, minimum height=9mm},
  sidenote/.style={font=\sffamily\scriptsize, text=octquiet},
  wire/.style={thick, octink},
]
  \foreach \k in {0,1,2,3} {
    \node[core] (c\k) at (\k*31mm, 0) {};
    \node[cpu] (cpu\k) at ($(c\k.north)+(0,-4.5mm)$) {CPU \k \quad OoO 8};
    \node[l1] (l1\k) at ($(c\k.south)+(0,4.5mm)$) {L1 \k \quad id \k};
    \draw[wire, <->] (cpu\k) -- (l1\k);
  }
  \node[font=\sffamily\scriptsize, text=octquiet, align=center, above=1mm of $(c1.north)!0.5!(c2.north)$]
    {\code{CacheControllerExclusive}, snoop MESI (\code{MESI\_splitBus\_snooping.csv}) \\ 16\,KB direct-mapped, 64\,B lines, untimed port, 16 MSHRs, 8-entry write-back buffer, queue 16};
  \node[bus, below=9mm of $(c1.south)!0.5!(c2.south)$, minimum width=120mm] (b0)
    {\textbf{bus[0]} \quad TripleBus, \code{SplitBusController}: request / response / service lanes \\ \scriptsize 256-deep buffers, 128 reserved for responses \textperiodcentered\ arbiter per lane: TDM by default (FCFS, RR selectable)};
  \foreach \k in {0,1,2,3} \draw[wire] (l1\k.south) -- (l1\k.south |- b0.north);
  \node[llc, below=8mm of b0, minimum width=120mm] (llc)
    {\textbf{LLC} \quad \code{CacheController\_End2End}, snoop MESI home (\code{MESI\_LLC.csv}) \quad id 10 \\ \scriptsize shared, inclusive \textperiodcentered\ 32\,KB, 2-way, 64\,B lines \textperiodcentered\ 10-cycle data port, FCFS port arbiter \textperiodcentered\ 64 MSHRs, 32-entry write-back buffer, queue 64};
  \draw[wire, <->] (b0.south) -- (llc.north);
  \node[bus, below=8mm of llc, minimum width=120mm] (b1)
    {\textbf{bus[1]} \quad \code{Point2Point}: LLC $\leftrightarrow$ memory, one message per direction per cycle};
  \draw[wire, <->] (llc.south) -- (b1.north);
  \node[mem, below=8mm of b1, minimum width=120mm] (mem)
    {\textbf{Main memory} \quad \code{MainMemoryController}, fixed 250-cycle latency, queue 32 \quad id 100 \\ \scriptsize or \code{MCsim} (DDR4, FR-FCFS) with \code{-p main\_memory\_type(s)=MCsim}};
  \draw[wire, <->] (b1.south) -- (mem.north);
  \node[sidenote, right=3mm of b0.east, align=left] {all clocks:\\period 100};
\end{tikzpicture}

\caption{The system \code{MultiCoreSystem.csv} builds. Four trace-driven cores with private
MESI L1s on a TripleBus, one shared inclusive LLC, a point-to-point memory bus and a
fixed-latency main memory. Every number is a configuration key (\cref{app:config-keys}).}
\label{fig:shipped-system}
\end{figure}

Running \code{Octopus\_Simulator -s MultiCoreSystem} instantiates the C++ system class
\code{MultiCoreSystem} and configures it from
\file{configuration/SystemConfigurations/MultiCoreSystem.csv}. \Cref{fig:shipped-system} is that
system; it is the one every route, datapath and sequence figure in this manual is drawn on, and
the one the tutorial runs first. The points that matter when reading a report:

\begin{itemize}
\item \textbf{Four cores} (\key{num\_cores}), each a trace-driven \code{CPU} with an
  out-of-order window of eight outstanding requests, and each with a private L1 of
  \code{CacheControllerExclusive} running snoop MESI --- 16\,KB, direct-mapped, 64-byte lines,
  an untimed data port, 16 MSHRs, an 8-entry write-back buffer and a 16-deep processing queue.
\item \textbf{\code{bus[0]}} is a TripleBus under a \code{SplitBusController}: three lanes ---
  request, response, service --- each with its own arbiter, TDM by default
  (\file{SplitBusController.csv}); 256-deep buffers with 128 entries reserved for in-flight
  responses so that the request lane stalls before responses can overflow.
\item \textbf{The LLC} (id 10) is a \code{CacheController\_End2End}: shared, inclusive, 32\,KB,
  2-way, 64-byte lines, a 10-cycle data port with an FCFS port arbiter, 64 MSHRs, a 32-entry
  write-back buffer and a 64-deep queue. Its \code{MESI\_LLC} FSM is the home of the protocol.
\item \textbf{\code{bus[1]}} is point-to-point between the LLC and memory (id 100).
\item \textbf{Main memory} is a \code{MainMemoryController} with a fixed 250-cycle latency;
  \code{-p "main\_memory\_type(s)=MCsim"} swaps in cycle-accurate DDR4 (\cref{ch:main-memory}).
\end{itemize}

Two sibling presets build the same topology with a different protocol:
\file{MultiCoreSystem_Snoop.csv} (snoop MSI) and \file{MultiCoreSystem_Directory.csv} (directory
MSI, pinned to an FCFS bus); \file{MultiCoreSystem_Mesh.csv} replaces the bus with a mesh
(\cref{ch:interconnects}). A protocol is five coupled keys, which is why these are presets and
not single \code{-p} overrides (\cref{ch:configuration}).

\begin{wherebox}
\file{src/ClockManager.cpp}, \file{header/ClockManager.h} --- the clock and
\code{ClockedObj}; \file{src/Configurable.cpp} --- parameters;
\file{header/Interconnect/CommunicationInterface.h} --- \code{Message} and the interfaces;
\file{src/SystemConfigurations/MultiCoreSystem.cpp} --- the system class that wires
\cref{fig:shipped-system}; \file{configuration/SystemConfigurations/MultiCoreSystem.csv} --- its
parameters; \file{Octopus_Simulator.cpp}, \file{src/CacheSim.cpp} --- \code{main} and the run.
\end{wherebox}
